\documentclass[12pt]{article}
\usepackage{amsmath, amssymb}
\usepackage[margin=1in]{geometry}
\setlength{\parskip}{6pt}

\title{QUBO Formulation: Credit Scorecard Combination Problem}
\date{}

\begin{document}
\maketitle

\subsection*{Credit Scorecard Combination Problem}

\paragraph{Problem Statement.}
Given a set of $N$ candidate scorecards, each associated with a non-negative expected revenue contribution $R_i$ for $i \in \{0, \dots, N-1\}$, and a symmetric non-negative pairwise overlap penalty $P_{ij}$ for $0 \leq i < j \leq N-1$, the objective is to select exactly $K$ scorecards such that the total expected revenue minus the sum of pairwise overlap penalties is maximized.

\paragraph{Decision Variables.}
For each candidate scorecard $i \in \{0, \dots, N-1\}$, introduce a binary decision variable $x_i \in \{0, 1\}$, where $x_i = 1$ indicates that scorecard $i$ is selected, and $x_i = 0$ indicates it is not selected.

\paragraph{QUBO Derivation.}
We transform the maximization problem into a minimization problem suitable for a Quadratic Unconstrained Binary Optimization (QUBO) formulation by negating the objective and incorporating the cardinality constraint via a penalty method.

\textbf{Step 1 -- Original Objective.} The original objective is to maximize revenue minus overlap penalties:
\begin{align}
    \max_{\mathbf{x}} \sum_{i=0}^{N-1} R_i x_i - \sum_{0 \leq i < j \leq N-1} P_{ij} x_i x_j.
\end{align}
For QUBO minimization, we equivalently minimize the negative of this expression:
\begin{align}
    \min_{\mathbf{x}} \left( -\sum_{i=0}^{N-1} R_i x_i + \sum_{0 \leq i < j \leq N-1} P_{ij} x_i x_j \right).
\end{align}

\textbf{Step 2 -- Constraint Formulation.} The requirement to select exactly $K$ scorecards is expressed as:
\begin{align}
    \sum_{i=0}^{N-1} x_i = K.
\end{align}

\textbf{Step 3 -- Penalty Term Expansion.} We enforce the constraint by adding a squared penalty term scaled by a positive weight $\lambda$:
\begin{align}
    \text{Penalty}(\mathbf{x}) = \lambda \left( \sum_{i=0}^{N-1} x_i - K \right)^2.
\end{align}
Expanding the square and applying the idempotent property $x_i^2 = x_i$ for binary variables yields:
\begin{align}
    \left( \sum_{i=0}^{N-1} x_i - K \right)^2 &= \sum_{i=0}^{N-1} x_i^2 - 2K \sum_{i=0}^{N-1} x_i + K^2 + \sum_{i \neq j} x_i x_j \\
    &= \sum_{i=0}^{N-1} x_i - 2K \sum_{i=0}^{N-1} x_i + K^2 + 2 \sum_{0 \leq i < j \leq N-1} x_i x_j \\
    &= (1 - 2K) \sum_{i=0}^{N-1} x_i + 2 \sum_{0 \leq i < j \leq N-1} x_i x_j + K^2.
\end{align}
Multiplying by $\lambda$ gives the general penalty contribution to the Hamiltonian:
\begin{align}
    \lambda \left[ (1 - 2K) \sum_{i=0}^{N-1} x_i + 2 \sum_{0 \leq i < j \leq N-1} x_i x_j + K^2 \right].
\end{align}

\textbf{Step 4 -- Combined Hamiltonian.} Adding the negated objective and the penalty term produces the full QUBO Hamiltonian $H(\mathbf{x})$:
\begin{align}
    H(\mathbf{x}) &= \left( -\sum_{i=0}^{N-1} R_i x_i + \sum_{0 \leq i < j \leq N-1} P_{ij} x_i x_j \right) + \lambda \left[ (1 - 2K) \sum_{i=0}^{N-1} x_i + 2 \sum_{0 \leq i < j \leq N-1} x_i x_j + K^2 \right].
\end{align}

\textbf{Step 5 -- Q Matrix Entries and Offset.} Collecting linear and quadratic terms yields the standard QUBO form $H(\mathbf{x}) = \sum_{i=0}^{N-1} Q_{i,i} x_i + \sum_{0 \leq i < j \leq N-1} Q_{i,j} x_i x_j + c$, with entries defined as:
\begin{align}
    Q_{i,i} &= -R_i + \lambda(1 - 2K), \quad \forall i \in \{0, \dots, N-1\}, \\
    Q_{i,j} &= P_{ij} + 2\lambda, \quad \forall 0 \leq i < j \leq N-1, \\
    c &= \lambda K^2.
\end{align}

\paragraph{Penalty Weight Justification.}
The penalty weight $\lambda$ must be chosen large enough to ensure that any violation of the cardinality constraint incurs a strictly higher energy cost than the maximum possible improvement in the objective function. The maximum objective gain achievable by selecting $K+1$ or $K-1$ items instead of $K$ is bounded by the sum of all individual revenues. Therefore, the condition
\begin{align}
    \lambda > \sum_{i=0}^{N-1} R_i
\end{align}
suffices to guarantee that infeasible solutions are strictly penalized relative to all feasible solutions.

\paragraph{Correctness Argument.}
Minimizing $H(\mathbf{x})$ over $\{0, 1\}^N$ recovers the optimal scorecard combination because (i) any assignment satisfying $\sum_{i=0}^{N-1} x_i = K$ incurs zero penalty, reducing $H(\mathbf{x})$ to the negated original objective; (ii) any infeasible assignment incurs a penalty of at least $\lambda$, which by construction exceeds the maximum possible reduction in the objective term achievable by any feasible assignment. Consequently, the global minimum of $H(\mathbf{x})$ must correspond to a feasible solution that minimizes the negated objective, which is equivalent to maximizing the original revenue-minus-penalties objective.

\paragraph{Illustrative Example.}
Consider a concrete instance with $N = 3$ candidate scorecards and $K = 2$ required selections. Let the revenues be $R_0 = 10$, $R_1 = 20$, and $R_2 = 30$. Let the pairwise overlap penalties be $P_{01} = 5$, $P_{02} = 8$, and $P_{12} = 12$. The sum of revenues is $10 + 20 + 30 = 60$, so we choose $\lambda = 61$ to satisfy the penalty weight condition.

Substituting these values into the general QUBO parameters yields:
\begin{align}
    Q_{0,0} &= -10 + 61(1 - 4) = -193, \\
    Q_{1,1} &= -20 + 61(1 - 4) = -203, \\
    Q_{2,2} &= -30 + 61(1 - 4) = -213, \\
    Q_{0,1} &= 5 + 2(61) = 127, \\
    Q_{0,2} &= 8 + 2(61) = 130, \\
    Q_{1,2} &= 12 + 2(61) = 134, \\
    c &= 61 \times 2^2 = 244.
\end{align}
The resulting Hamiltonian is:
\begin{align}
    H(x_0, x_1, x_2) = -193x_0 - 203x_1 - 213x_2 + 127x_0x_1 + 130x_0x_2 + 134x_1x_2 + 244.
\end{align}
Enumerating the three feasible assignments ($\sum x_i = 2$):
\begin{align}
    H(1, 1, 0) &= -193 - 203 + 127 + 244 = -25, \\
    H(1, 0, 1) &= -193 - 213 + 130 + 244 = -32, \\
    H(0, 1, 1) &= -203 - 213 + 134 + 244 = -38.
\end{align}
The global minimum is $-38$ at $\mathbf{x}^* = (0, 1, 1)$. This corresponds to selecting scorecards $1$ and $2$, yielding a maximum total score of $R_1 + R_2 - P_{12} = 20 + 30 - 12 = 38$, which matches the negated minimum energy $-(-38) = 38$.
\end{document}